\documentclass{article}
\usepackage{amsmath, amssymb, amsfonts}

\newcommand{\op}{\prime}
\newcommand{\pdt}[2]{\frac{\partial{#1}}{\partial{#2}}}
\newcommand{\spdt}[2]{\frac{\partial^2{#1}}{\partial{#2}^2}}
\newcommand{\spdtm}[3]{\frac{\partial^2{#1}}{\partial{#2}\partial{#3}}}
\DeclareMathOperator{\tr}{Tr}
\title{Sound Waves at Low Frequency}
\begin{document}
\author{Amey Joshi}
\maketitle
\section{Low-frequency regime}\label{s1}
The inter-atomic distances and lattice constants of crystals are in the range
of a few \AA. Waves whose wavelengths are much longer than these distances are
insensitive to changes at the lattice levels. We will focus on waves whose 
wavelengths are longer than $10^{-6}$ cm. The speed of sound in solids is in a
few thousand metres per second. Therefore, we shall consider waves with a 
frequency less than $10^{11}$ Hz. We shall call these the low frequency waves.
Note that this is quite different from the range of sound waves considered to be
of low frequency in acoustics and electronics. In those fields, the ability of
humans to hear waves is taken into account.

In the low-frequency regime a crystal can be considered as a continuum and we
can use the classical theory of elasticity to analyse the propagation of sound
waves.

\section{The strain tensor}\label{s2}
For notational convenience, we shall label the three cartesian axes as $x_1, 
x_2, x_3$ instead of $x, y, z$ and the components of a vector $\vec{F}$ shall
be denoted as $F_1, F_2, F_3$ instead of $F_x, F_y, F_z$.

We shall start with the definition of deformation of a solid. Fix the cartesian
axis in an unstrained solid. After deformation, the axes $x_i$ are transformed
to axes $x_i^\op$. Unit vectors along these axes $\hat{x}_i$ are transformed to
$\vec{x^\op}_i$. Note that a unit vector may be transformed into a vector whose
magnitude need not be $1$. We define numbers $\epsilon_{ij}$ by
\begin{equation}\label{e1}
\vec{x^\op}_i - \hat{x}_i = \epsilon_{i1}\hat{x}_1 + \epsilon_{i2}\hat{x}_2
+ \epsilon_{i3}\hat{x}_3 = \epsilon_{ij}\hat{x}_j,
\end{equation}
where we have used the summation convention. We can express \eqref{e1} as
\begin{equation}\label{e2}
\vec{x^\op}_i = (\delta_{ij} + \epsilon_{ij})\hat{x}_j,
\end{equation}
where $\delta_{ij}$ is the Kronecker-$\delta$. Up to first order in 
$\epsilon_{ij}$, 
\begin{equation}\label{e3}
\lVert\vec{x^\op}_i\rVert^2 = 1 + 2\epsilon_{ii},
\end{equation}
from which it follows that
\begin{equation}\label{e4}
\lVert\vec{x^\op}_i\rVert = 1 + \epsilon_{ii} = \lVert\hat{x}_i\rVert +
\epsilon_{ii}.
\end{equation}
Thus the change in length of unit vectors is
\begin{equation}\label{e5}
\lVert\vec{x^\op}_i\rVert - \lVert\hat{x}_i\rVert = \epsilon_{ii}.
\end{equation}
The angles between the deformed axes are no longer equal to $\pi/2$. Instead,
they are
\begin{eqnarray*}
\cos\theta_{ij} &=& \frac{\vec{x^\op}_i \cdot \vec{x^\op}_j}
{\lVert\vec{x^\op}_i\rVert\lVert\vec{x^\op}_j\rVert}  \\
&=& \frac{(\delta_{il} + \epsilon_{il})(\delta_{jm} + \epsilon_{jm})\hat{x}_l
\cdot\hat{x}_m}{(1 + \epsilon_{ii})(1 + \epsilon_{jj})}  \\
&=& \frac{(\delta_{il} + \epsilon_{il})(\delta_{jm} + \epsilon_{jm})\delta_{lm}}
{(1 + \epsilon_{ii} + \epsilon_{jj})}  \\
&=& \frac{(\delta_{il} + \epsilon_{il})(\delta_{jl} + \epsilon_{jl})}
{(1 + \epsilon_{ii} + \epsilon_{jj})}  \\
&=& \frac{\delta_{ij} + \epsilon_{ij} + \epsilon_{ji}}
{(1 + \epsilon_{ii} + \epsilon_{jj})}  
\end{eqnarray*}
If $i = j$ then $\cos\theta_{ii} = 1$ so that $\theta_{ii} = 0$, as we would 
have expected. When $i \ne j$, up to first order in $\epsilon_{ij}$, we have
\begin{equation}\label{e6}
\cos\theta_{ij} = \epsilon_{ij} + \epsilon_{ji}.
\end{equation}
We define the strain tensor, $e_{ab}$ in abstract index notation, in terms of 
$\epsilon_{ij}$ as
\begin{equation}\label{e7}
e_{ab} = \begin{bmatrix}
\epsilon_{11} & \epsilon_{12} + \epsilon_{21} & \epsilon_{13} + \epsilon_{31} \\
\epsilon_{21} + \epsilon_{12} & \epsilon_{22} & \epsilon_{13} + \epsilon_{31} \\
\epsilon_{31} + \epsilon_{13} & \epsilon_{32} + \epsilon_{23} & \epsilon_{33} 
\end{bmatrix}.
\end{equation}
Deformation also causes the volume of a parallelepiped formed by $\hat{x}_i$
to change from $1$ to $|\vec{x^\op}_1\cdot(\vec{x^\op}_2\times\vec{x^\op}_3)|$,
which is
\[
\begin{vmatrix}
1 + \epsilon_{11} & \epsilon_{12} & \epsilon_{13} \\
\epsilon_{21} & 1 + \epsilon_{22} & \epsilon_{23} \\
\epsilon_{31} & \epsilon_{32} & 1 + \epsilon_{33}
\end{vmatrix} = 1 + \epsilon_{11} + \epsilon_{22} + \epsilon_{33} = 
1 + \tr{e_{aa}}.
\]
We define the dilation of the volume element, $\delta$ as the ratio of the 
change in volume to the original volume. Thus, the dilation
\begin{equation}\label{e8}
\delta = \tr{e_{aa}} = e_{11} + e_{22} + e_{33},
\end{equation}
is equal to the trace of the strain tensor when the deformation is slight.

\section{Stress tensor}
Consider a plane portion with unit area of a solid aligned such that the $x$-axis 
is perpendicular to it. A force acting on this surface will, in general, have all
three components. We denote them as $(T_{xx}, T_{xy}, T_{xz})$. Forces acting
plane surfaces perpendicular to other axes are denoted analogously. Thus
$T_{ij}$ is the $j$-th component of the force acing per unit area aligned along
the $i$-axis.

Consider a small volume element such that the forces on opposite faces are 
equal. If $T_{xy}$ and $T_{yx}$ are not equal then there will be a net torque
on the element resulting in an angular motion about the $z$-axis. If the volume
element is in equilibrium, then the two forces have to be equal. This is true
for all combinations of unequal indices. That is
\begin{equation}\label{e9}
T_{ij} = T_{ji}.
\end{equation}

Hooke's law states that the stress and strain are proportional to each other.
Thus,
\begin{eqnarray}
e_{ab} &=& S_{abcd}T_{cd}  \label{e10} \\
T_{ab} &=& C_{acbd}e_{cd}, \label{e11}
\end{eqnarray}
where $S_{abcd}$ is called the elastic compliance tensor and $C_{abcd}$ is 
called the elastic stiffness tensor, also called the modulus of elasticity.
From equations \eqref{e7} and \eqref{e9} it is evident that the strain and
stress tensors are both symmetric. Each one of them has $6$ independent 
components. Each of these components results in a corresponding element in 
$C_{abcd}$. Thus, there are $36$ independent components in $C_{abcd}$. For an
elastic material deforming reversibly the work done is stored as elastic energy.
\[
U = \frac{1}{2} C_{abcd}e_{ab}e_{cd}.
\]
Since
\[
\frac{\partial^2 U}{\partial e_{ab}\partial e_{cd}} = \frac{\partial^2 U}
{\partial e_{cd}\partial e_{ab}}
\]
we must also have $C_{abcd} = C_{cdab}$. We can represent the double index
$ab$ by $\alpha$ and $cd$ by $\beta$ and treat $C_{abcd}$ as if it were a 
$6 \times 6$ matrix. The symmetric $C_{\alpha\beta} = C_{\beta\alpha}$ reduces
the number of independent components to $21$. (It is the sum of $6$ diagonal
elements plus $5 + 4 + 3 + 2 + 1 = 15$ off-diagonal elements.) The same is true
for the elastic compliance tensor.

\subsection{Alternative notation}
The six independent components of the strain and stress tensor can be chosen to
be $xx, yy, zz, yz, zx, xy$. They can be relabelled as $1, \ldots, 6$. We can
treat them to be six dimension vectors $e_\mu$ and $T_\mu$. We can write 
equations \eqref{e10} and \eqref{e11} as
\begin{eqnarray}
e_\mu &=& S_{\mu\nu}T_\nu \label{e12} \\
T_\mu &=& C_{\mu\nu}T_\nu \label{e13}
\end{eqnarray}
where $S_{\mu\nu}$ and $C_{\mu\nu}$ are $6 \times 6$ symmetric compliance and 
stiffness matrices.

\section{Elastic waves in cubic crystals}
We start with Kittel's expression of \eqref{e13} for a cubic crystal,
\begin{equation}\label{e14}
\begin{pmatrix}T_1 \\ T_2 \\ T_3 \\ T_4 \\ T_5 \\ T_6\end{pmatrix} = 
\begin{pmatrix}
C_{11} & C_{12} & C_{12} & 0 & 0 & 0 \\
C_{12} & C_{11} & C_{12} & 0 & 0 & 0 \\
C_{12} & C_{12} & C_{11} & 0 & 0 & 0 \\
0 & 0 & 0 & C_{44} & 0 & 0 \\
0 & 0 & 0 & 0 & C_{44} & 0 \\
0 & 0 & 0 & 0 & 0 & C_{44}
\end{pmatrix}
\begin{pmatrix}e_1 \\ e_2 \\ e_3 \\ e_4 \\ e_5 \\ e_6\end{pmatrix}
\end{equation}
In order to write the equations of motion, it is convenient to use the 
standard cartesian tensor notation for the components of stress and strain.
However, we keep the notation for the components of the stiffness tensor 
unchanged. Thus, \eqref{e14} becomes
\begin{equation}\label{e15}
\begin{pmatrix}T_{xx} \\ T_{yy} \\ T_{zz} \\ T_{yz} \\ T_{zx} \\ T_{xy}
\end{pmatrix} = 
\begin{pmatrix}
C_{11} & C_{12} & C_{12} & 0 & 0 & 0 \\
C_{12} & C_{11} & C_{12} & 0 & 0 & 0 \\
C_{12} & C_{12} & C_{11} & 0 & 0 & 0 \\
0 & 0 & 0 & C_{44} & 0 & 0 \\
0 & 0 & 0 & 0 & C_{44} & 0 \\
0 & 0 & 0 & 0 & 0 & C_{44}
\end{pmatrix}
\begin{pmatrix}e_{xx} \\ e_{yy} \\ e_{zz} \\ e_{yz} \\ e_{zx} \\ e_{xy}
\end{pmatrix}
\end{equation}
The Cauchy equation for continua in abstract index notation is
\begin{equation}\label{e16}
\rho\frac{\partial^2 R_b}{\partial t^2} = \nabla_{a}T_{ab},
\end{equation}
where 
\begin{equation}\label{e17}
\vec{R}(\vec{r}) = \vec{r^\op} - \vec{r} = u_i(\vec{r})\hat{x}_i
\end{equation}
is the deformation. Now,
\[
\nabla_a T_{ab} = \hat{e}_i\frac{\partial}{\partial x_i}\cdot\hat{e}_j\hat{e}_k
T_{jk} = \delta_{ij}\frac{\partial}{\partial x_i}\hat{e}_k T_{jk} = 
\frac{\partial T_{jk}}{\partial x_j}\hat{e}_k = 
\frac{\partial T_{kj}}{\partial x_j}\hat{e}_k
\]
where the last equations follows from the symmetry of the stress tensor. Thus,
\begin{eqnarray*}
\nabla_a T_{ab} &=& \hat{e}_x\left(\frac{\partial T_{xx}}{\partial x} + 
\frac{\partial T_{xy}}{\partial y}+\frac{\partial T_{xz}}{\partial z}\right)+\\
 & & \hat{e}_y\left(\frac{\partial T_{yx}}{\partial x} + 
\frac{\partial T_{yy}}{\partial y}+\frac{\partial T_{yz}}{\partial z}\right)+\\
 & & \hat{e}_z\left(\frac{\partial T_{zx}}{\partial x} + 
\frac{\partial T_{zy}}{\partial y} + \frac{\partial T_{zz}}{\partial z}\right)
\end{eqnarray*}
Using \eqref{e15}
\begin{eqnarray*}
\nabla_a T_{ab} &=& \hat{e}_x\left(C_{11}\pdt{e_{xx}}{x} + C_{12}\pdt{e_{yy}}{x}
+C_{12}\pdt{e_{zz}}{x}+C_{44}\pdt{e_{xy}}{y} + C_{44}\pdt{e_{zx}}{x}\right)+\\
& & \hat{e}_y\left(C_{44}\pdt{e_{xy}}{x} + C_{12}\pdt{e_{xx}}{y} + 
C_{11}\pdt{e_{yy}}{y} + C_{12}\pdt{e_{zz}}{y} + C_{44}\pdt{e_{yz}}{z}\right)+\\
& & \hat{e}_z\left(C_{44}\pdt{e_{zx}}{x} + C_{44}\pdt{e_{yz}}{y} +
C_{12}\pdt{e_{xx}}{z} + C_{12}\pdt{e_{yy}}{z} + C_{11}\pdt{e_{zz}}{z}\right)
\end{eqnarray*}
Using these expressions in \eqref{e16},
\begin{eqnarray*}
\rho\spdt{u_x}{t} &=& C_{11}\pdt{e_{xx}}{x} + C_{12}\left(\pdt{e_{yy}}{x} +
\pdt{e_{zz}}{x}\right) + C_{44}\left(\pdt{e_{xy}}{y} + \pdt{e_{zx}}{x}\right) \\
\rho\spdt{u_y}{t} &=& C_{11}\pdt{e_{yy}}{y} + C_{12}\left(\pdt{e_{xx}}{y} +
\pdt{e_{zz}}{y}\right) + C_{44}\left(\pdt{e_{xy}}{x} + \pdt{e_{yz}}{z}\right) \\
\rho\spdt{u_z}{t} &=& C_{11}\pdt{e_{zz}}{z} + C_{12}\left(\pdt{e_{xx}}{z} +
\pdt{e_{yy}}{z}\right) + C_{44}\left(\pdt{e_{zx}}{x} + \pdt{e_{yz}}{y}\right)
\end{eqnarray*}
The strains are derivatives of $u_x, u_y, u_z$ so that these equations become
\begin{eqnarray*}
\rho\spdt{u_x}{t} &=& C_{11}\spdt{u_x}{x} + C_{12}\left(\spdtm{u_y}{x}{y} +
\spdtm{u_z}{x}{z}\right) + \\
& & C_{44}\left(\spdt{u_x}{y} + \spdtm{u_y}{x}{y} + 
\spdtm{u_x}{x}{z} + \spdt{u_x}{z}\right) \\
\rho\spdt{u_y}{t} &=& C_{11}\spdt{u_y}{y} + C_{12}\left(\spdtm{u_x}{x}{y} +
\spdtm{u_z}{y}{z}\right) + \\
 & & C_{44}\left(\spdtm{u_x}{x}{y} + \spdt{u_y}{x} + \spdt{u_y}{z}
+ \spdtm{u_z}{y}{z}\right) \\
\rho\spdt{u_z}{t} &=& C_{11}\pdt{u_z}{z} + C_{12}\left(\spdtm{u_x}{x}{z} +
\pdt{u_y}{y}{z}\right) + \\
& & C_{44}\left(\spdtm{u_z}{x}{z} + \spdt{u_z}{x} + \spdtm{u_y}{y}{z}
+ \spdt{u_z}{y}\right)
\end{eqnarray*}
Simplifying them a bit, we get the equations of elastic waves in a cubic crystal
\begin{eqnarray}
\rho\spdt{u_x}{t} &=& C_{11}\spdt{u_x}{x} +C_{44}\left(\spdt{u_x}{y} + 
\spdt{u_x}{z}\right) + \nonumber \\
& & (C_{12} + C_{44})\left(\spdtm{u_y}{x}{y} + \spdtm{u_z}{x}{z}\right)
\label{e18} \\
\rho\spdt{u_y}{t} &=& C_{11}\spdt{u_y}{y} + C_{44}\left(\spdt{u_y}{x} + 
\spdt{u_y}{z}\right) + \nonumber \\
& & (C_{12} + C_{44})\left(\spdtm{u_x}{x}{y} + \spdtm{u_z}{y}{z}\right)
\label{e19} \\
\rho\spdt{u_z}{t} &=& C_{11}\spdt{u_z}{z} + C_{44}\left(\spdt{u_z}{x} + 
\spdt{u_z}{z}\right) + \nonumber \\
& & (C_{12} + C_{44})\left(\spdtm{u_x}{x}{z} + \spdtm{u_y}{y}{z}\right)
\label{e20}
\end{eqnarray}
\end{document}
